\begin{lemma}[Common exterior operators and native boundary transport]%
    \label{lem:peps_regular_bond_extension_transport}
    \lean{TNLean.PEPS.regularTwistedSite_regularRegionBondExtension_eq_outside,
        TNLean.PEPS.openRegionWeight_regularTreeCycleBondExtension}
    \leanok
    \uses{thm:peps_regular_region_gauge_contraction}
    Prescribe internal bond operators in a finite region $R$ and retain one
    common assignment $u$ on exterior and crossing bonds. Altering the internal
    operators leaves every exterior site tensor unchanged. If the original
    sites are invariant under simultaneous regular translation, every
    tree-cycle assignment $\omega$ satisfies
    \begin{align}
        A_R(\omega;u;\theta)
        =A_R(\omega;1;\mathcal B_u\theta),
    \end{align}
    where $\mathcal B_u$ is the native boundary transport for the identity
    vertex gauge. It depends on $u$ but not on $\omega$. Here a graph on $R$
    specifies which internal bonds carry the cycle assignment; a spanning-tree
    hypothesis is unnecessary for these identities. These are the shared
    contraction facts in the operations of
    \cite[Theorems~6.16--6.17]{Schuch2010PEPS}.
\end{lemma}

\begin{proof}\leanok
    An edge incident to an exterior vertex cannot be wholly internal to $R$,
    so its operator is the common exterior one. A crossing edge is likewise
    not internal. The identity vertex gauge therefore retains each prescribed
    internal operator and transfers the common crossing operators into exactly
    the same boundary transport for every internal assignment. The actual
    gauge-contraction identity gives the displayed open-region equality.
\end{proof}

\begin{theorem}[Global contraction of the two-cycle flux operation]%
    \label{thm:peps_regular_two_cycle_global_flux_move}
    \lean{TNLean.PEPS.regularRegionSiteExtension,
        TNLean.PEPS.regionLocalTerm_mem_unitaryGroup,
        TNLean.PEPS.exists_unitary_regularTwoCycleGlobalFluxMove,
        TNLean.PEPS.regularRegionBondExtension,
        TNLean.PEPS.exists_unitary_regularTwoCycleGlobalFluxMove_bondOperators}
    \leanok
    \uses{lem:peps_regular_bond_extension_transport,
        thm:peps_regular_two_cycle_physical_flux_move,
        thm:peps_open_region_cut,
        thm:peps_regular_region_gauge_contraction}
    Let the original regular site maps be $G$-isometric. Choose a spanning
    tree and a root of a finite region $R$, together with distinct internal
    bonds $e_0,e_1$ outside the tree. Let $u_g$ have value $g$ on $e_0$ and
    identity on every other internal bond; let $u'_g$ have value $g$ on
    $e_0,e_1$ and identity on every other internal bond.
    There is one physical unitary $W$ on $R$ whose extension
    $\widehat W=W\otimes I_{R^c}$ is unitary on the whole system and satisfies
    \begin{align}
        \widehat W\Psi(u_g,b)=\Psi(u'_g,b)
    \end{align}
    for every $g\in G$ and every common family $b$ of complementary site
    coefficients. Here both vectors are the actual finite graph contractions,
    using the inserted original sites in $R$ and $b$ outside $R$.

    There is also such a fixed identity-extended unitary with the following
    bond-operator formulation. For an arbitrary common assignment $u$ on
    exterior and crossing bonds, retain $u$ on those bonds and prescribe
    $u_g$, respectively $u'_g$, on internal bonds. Denote the resulting
    assignments by $\widetilde u_g$ and $\widetilde u'_g$. Then
    \begin{align}
        \widehat W\Psi(\widetilde u_g)
        =\Psi(\widetilde u'_g)
        \qquad (g\in G).
    \end{align}
    The unitary is chosen before both $g$ and the common exterior data.
    Crossing operators need not be trivial, and no contraction or Gram
    identity is a hypothesis. Linearity therefore permits coherent
    superpositions of flux labels.
    This is the global contraction consequence of the auxiliary two-cycle
    operation for \cite[Theorem~6.16]{Schuch2010PEPS}. Its six-vertex torus
    instance is established below. The unrestricted string statement and
    parent-Hamiltonian conclusions remain separate.
\end{theorem}

\begin{proof}\leanok
    Extend the regional unitary by the complementary identity and transport
    it through the physical configuration partition. Kronecker products
    and reindexing preserve unitarity. For fixed complementary spins, its
    action is exactly the regional matrix action.
    The actual cut formula gives
    \begin{align}
        \Psi(\sigma,\tau)
        =\sum_{\theta}\Psi_R(\sigma,\theta)
        \Psi_{R^c}(\tau,\theta).
    \end{align}
    Apply the all-boundary-column identity to the first factor. The second
    factor is arbitrary and common, so the sum gives the global equality.
    For the bond-operator formulation, the identity vertex gauge moves
    each crossing operator into its label at the region endpoint. This
    boundary transport is the same for the one- and two-insertion families;
    their internal operators are exactly the prescribed tree and cycle
    assignments. Their complementary sites also agree, because the changed
    bonds are wholly inside $R$. The same cut argument proves the equality
    for every common exterior and crossing assignment.
\end{proof}
